Multiplying representatives by units does not change whether $x/y\in A$, so the order is well-defined. Reflexivity and transitivity follow from $1\in A$ and closure under multiplication. If both $x/y$ and $y/x$ lie in $A$, then $x/y\in U$, giving $[x]=[y]$. The valuation-ring property gives totality. Finally, $(xz)/(yz)=x/y$, so multiplication by $[z]$ preserves the order.

Suppose $v(x)\geq v(y)$. Then $x/y\in A$, so $(x+y)/y=1+x/y\in A$. Hence $v(x+y)\geq v(y)=\min(v(x),v(y))$.
